\chapter{Programs — Logic}
%%% generated by the blueprint pipeline from TCSlib.Programs.Light.Logic
\section{Overview}

Declarations of the Lean module \texttt{TCSlib.Programs.Light.Logic}.

\section{Declarations}

\begin{theorem}[Preserving an indexed entry under concatenation]
\lean{Light.getElem?_append_of_eq_some}
\leanok
\uses{Light.Program, Light.Stmt}
\difficulty{3}
Let $P$ and $R$ be finite lists of statements, let $p$ be a natural number, and let $b$
be a statement. If the entry at index $p$ of $P$ is $b$, then the entry at index $p$ of
$P\mathbin{++}R$ is also $b$.
\end{theorem}

\begin{theorem}[Indexing a concatenated program]
\lean{Light.getElem?_append_append}
\leanok
\uses{Light.Program, Light.Stmt, Light.getElem?_append_of_eq_some}
\difficulty{3}
Let $P_0$, $B$, and $R$ be finite lists of statements, and let $i$ be a natural number.
If the entry at index $i$ of $B$ is a statement $b$, then the entry at index $|P_0|+i$
of the concatenated list $P_0\mathbin{++}B\mathbin{++}R$ is $b$.
\end{theorem}

\begin{theorem}[Indexing a program after concatenation]
\lean{Light.getElem?_append_length_add}
\leanok
\uses{Light.Program, Light.Stmt}
\difficulty{3}
Let $P_0$ and $R$ be finite lists of statements, and let $i$ be a natural number. If the
entry at index $i$ of $R$ is a statement $b$, then the entry at index $|P_0|+i$ of
$P_0\mathbin{++}R$ is $b$.
\end{theorem}

\begin{theorem}[Execution under an extended program]
\lean{Light.Exec.append}
\label{Light.Exec.append}
\leanok
\uses{Light.Exec, Light.Program, Light.State, Light.Stmt, Light.getElem?_append_of_eq_some}
\difficulty{4}
Fix limits, a program $P$, a depth $d$, a statement $s$, states $\sigma$ and $\sigma'$,
and a cost $c$. If $s$ executes from $\sigma$ to $\sigma'$ at depth $d$ under $P$ with
cost $c$, then for every program $R$, it also executes with the same states, depth, and
cost under the concatenation of $P$ and $R$.
\end{theorem}

\begin{definition}[Execution within a cost bound]
\lean{Light.Ends}
\label{Light.Ends}
\leanok
\uses{Light.Exec, Light.Limits, Light.Program, Light.State, Light.Stmt}
Given limits, a program, a depth parameter, a statement, an initial state, a
natural-number cost bound $T$, and a predicate $Q$ on states, the execution satisfies
this condition exactly when there are a final state and a cost such that executing the
statement from the initial state reaches that final state with that cost, the cost is at
most $T$, and $Q$ holds of the final state.
\end{definition}

\begin{theorem}[Monotonicity of Bounded Execution]
\lean{Light.Ends.mono}
\label{Light.Ends.mono}
\leanok
\uses{Light.Ends}
\difficulty{4}
Let $\mathrm{Ends}(s,\sigma,T,Q)$ mean that statement $s$ has an execution from state
$\sigma$ to some state satisfying $Q$, with cost at most $T$. For any limits, program,
depth, statement, initial state, natural numbers $T,T'$, and state predicates $Q,Q'$, if
$\mathrm{Ends}(s,\sigma,T,Q)$, $T\le T'$, and every state satisfying $Q$ also satisfies
$Q'$, then $\mathrm{Ends}(s,\sigma,T',Q')$.
\end{theorem}

\begin{theorem}[Execution under program extension]
\lean{Light.Ends.append}
\label{Light.Ends.append}
\leanok
\uses{Light.Ends, Light.Exec.append, Light.Program}
\difficulty{4}
Let $L$ be a set of execution limits, let $P$ and $R$ be programs, let
$d,T\in\mathbb{N}$, let $s$ be a statement, let $\sigma$ be an initial state, and let
$Q$ be a predicate on states. If $s$ has an execution from $\sigma$ under $P$ at depth
$d$ that ends in a state satisfying $Q$ and has cost at most $T$, then $s$ has an
execution from $\sigma$ under the concatenated program $P++R$ at depth $d$ that ends in
a state satisfying $Q$ and has cost at most $T$.
\end{theorem}

\begin{theorem}[Skip statement]
\lean{Light.Ends.skip}
\label{Light.Ends.skip}
\leanok
\uses{Light.Ends, Light.State}
\difficulty{1}
For any limit record, program, execution depth, state $\sigma$, budget $T$, and
postcondition $Q$, if $Q(\sigma)$ holds, then the skip statement can execute from
$\sigma$ with cost at most $T$ and end in a state satisfying $Q$.
\end{theorem}

\begin{theorem}[Bounded execution of an assignment]
\lean{Light.Ends.set}
\label{Light.Ends.set}
\leanok
\uses{Light.Ends, Light.Expr.Safe, Light.Expr.cost, Light.Expr.val, Light.State}
\difficulty{1}
Let $\mathrm{lim}$ be resource limits, $P$ a program, $d$ a call depth, $x$ a local index,
state, and procedure index $p$ whose body in $P$ is $body$. Let $args$ be a list of argument
$\sigma$, $\mathrm{cost}(e)+1\le T$, and $Q$ holds in the state obtained by updating the
local map of $\sigma$ at $x$ to $e$ evaluated in $\sigma$, then assigning $e$ to $x$
from $\sigma$ ends within time $T$ in a state satisfying $Q$.
\end{theorem}

\begin{theorem}[Store statement]
\lean{Light.Ends.store}
\label{Light.Ends.store}
\leanok
\uses{Light.Ends, Light.Expr.Safe, Light.Expr.cost, Light.Expr.val, Light.Limits.Addr, Light.State}
\difficulty{1}
Let a limit record, program, execution depth, state $\sigma$, budget $T$, address
expression $a$, value expression $e$, and postcondition $Q$ be given. Suppose $a$ and
$e$ are safe in $\sigma$, the value of $a$ is a valid address (nonnegative and less than
the space bound), and $\mathrm{cost}(a)+\mathrm{cost}(e)+1\le T$. If $Q$ holds after
updating memory at the natural-number address corresponding to the value of $a$ to the
value of $e$, while leaving the local function unchanged, then the store statement can
execute from $\sigma$ with cost at most $T$ and end in a state satisfying $Q$.
\end{theorem}

\begin{theorem}[Sequential composition of bounded executions]
\lean{Light.Ends.seq}
\label{Light.Ends.seq}
\leanok
\uses{Light.Ends, Light.State}
\difficulty{4}
Let $\mathrm{lim}$ be resource limits, $P$ a program, $d$ a call depth, and $s_1,s_2$
statements. If $s_1$ executes from state $\sigma$ within time $T_1$ to a state from which
$s_2$ executes within time $T_2$ and ends in a state satisfying $Q$, and $T_1+T_2\le T$,
then the sequential composition $s_1;s_2$ executes from $\sigma$ within time $T$ and ends
in a state satisfying $Q$.
\end{theorem}

\begin{theorem}[Bounded execution of a conditional]
\lean{Light.Ends.ite}
\label{Light.Ends.ite}
\leanok
\uses{Light.Cond.Holds, Light.Cond.Safe, Light.Cond.cost, Light.Ends, Light.State}
\difficulty{4}
Fix a program environment, resource limits, and execution depth. Let $c$ be a safe
condition in an initial state $\sigma$, and let $s_1,s_2$ be statements. Suppose that if
$c$ holds, then $s_1$ can execute from $\sigma$ with cost at most $T'$ to a state
satisfying $Q$, and if $c$ does not hold, then $s_2$ can do so. If
$\mathrm{cost}(c)+1+T'\le T$, the conditional statement executes from $\sigma$ with cost
at most $T$ to a state satisfying $Q$.
\end{theorem}

\begin{theorem}[One bounded loop iteration]
\lean{Light.Ends.whileStep}
\label{Light.Ends.whileStep}
\leanok
\uses{Light.Cond.Holds, Light.Cond.Safe, Light.Cond.cost, Light.Ends, Light.State}
\difficulty{4}
Let a program, depth bound, condition $c$, body statement $s$, initial state $\sigma$,
time bounds $T_1,T_2,T\in\mathbb N$, and postcondition $Q$ be given. If $c$ is safe and
true at $\sigma$, and the body can execute from $\sigma$ within time $T_1$ to a state
from which the same loop can end in a state satisfying $Q$ within time $T_2$, and
$\operatorname{cost}(c)+1+T_1+T_2\le T$, then the loop ends from $\sigma$ in a state
satisfying $Q$ within time $T$.
\end{theorem}

\begin{theorem}[Loop terminates immediately]
\lean{Light.Ends.whileDone}
\label{Light.Ends.whileDone}
\leanok
\uses{Light.Cond.Holds, Light.Cond.Safe, Light.Cond.cost, Light.Ends, Light.State}
\difficulty{1}
Let a program, depth bound, condition $c$, body statement $s$, initial state $\sigma$,
time bound $T\in\mathbb N$, and postcondition $Q$ be given. If $c$ is safe and false at
$\sigma$, $Q$ holds at $\sigma$, and $\operatorname{cost}(c)+1\le T$, then the loop with
condition $c$ and body $s$ ends in a state satisfying $Q$ within time $T$.
\end{theorem}

\begin{theorem}[Bounded loop from an invariant]
\lean{Light.Ends.while}
\label{Light.Ends.while}
\leanok
\uses{Light.Cond.Holds, Light.Cond.Safe, Light.Cond.cost, Light.Ends, Light.State}
\difficulty{6}
Let a program, depth bound, condition $c$, body $s$, initial state $\sigma$,
postcondition $Q$, invariant predicates $I_i$ on states for natural numbers $i$,
iteration count $n\in\mathbb N$, and per-iteration bounds $b_i\in\mathbb N$ be given. If
$I_0$ holds at $\sigma$, then for every $i<n$, every state satisfying $I_i$ has a safe
condition that holds and admits an execution of $s$ within time $b_i$ to a state
satisfying $I_{i+1}$. If every state satisfying $I_n$ has a safe condition that does not
hold and satisfies $Q$, then the loop with condition $c$ and body $s$ ends in a state
satisfying $Q$ within time
$\sum_{i=0}^{n-1}\bigl((\operatorname{cost}(c)+1+b_i)+(\operatorname{cost}(c)+1)\bigr)$.
\end{theorem}

\begin{theorem}[Bounded loop with fixed iteration cost]
\lean{Light.Ends.whileConst}
\label{Light.Ends.whileConst}
\leanok
\uses{Light.Cond.Holds, Light.Cond.Safe, Light.Cond.cost, Light.Ends, Light.Ends.mono, Light.Ends.while, Light.State}
\difficulty{1}
Let a program, depth bound, condition $c$, body statement $s$, initial state $\sigma$,
time bound $T\in\mathbb N$, postcondition $Q$, invariant predicates $I_i$ on states for
natural numbers $i$, and natural numbers $n,b$ be given. If $I_0$ holds at $\sigma$,
then for every $i<n$, every state satisfying $I_i$ has a safe condition that holds and
admits an execution of $s$ within time $b$ to a state satisfying $I_{i+1}$. If every
state satisfying $I_n$ has a safe condition that does not hold and satisfies $Q$, and
$n(\operatorname{cost}(c)+1+b)+\operatorname{cost}(c)+1\le T$, then the loop with
condition $c$ and body $s$ ends in a state satisfying $Q$ within time $T$.
\end{theorem}

\begin{theorem}[Termination of a while loop by a variant]
\lean{Light.Ends.whileVariant}
\label{Light.Ends.whileVariant}
\leanok
\uses{Light.Cond.Holds, Light.Cond.Safe, Light.Cond.cost, Light.Ends, Light.Ends.mono, Light.State}
\difficulty{6}
Let a program execute under limits and at recursion depth $d$, and consider a while loop
with guard $c$ and body $s$, starting in state $\sigma$. Suppose an invariant $I$ holds
initially, and whenever $I$ holds the guard is safe. Whenever $I$ holds and the guard is
true, the body terminates within cost $b$, preserves $I$, and strictly decreases a
natural-valued measure $m$. Whenever $I$ holds and the guard is false, a postcondition
$Q$ holds. If $m(\sigma)(\operatorname{cost}(c)+1+b)+\operatorname{cost}(c)+1\le T$,
then the loop terminates within cost $T$ and its final state satisfies $Q$.
\end{theorem}

\begin{theorem}[Cost of a procedure call]
\lean{Light.Ends.call}
\label{Light.Ends.call}
\leanok
\uses{Light.Ends, Light.Expr.Safe, Light.Expr.cost, Light.Expr.val, Light.State, Light.frame}
\difficulty{2}
Let $\mathrm{lim}$ be resource limits, $P$ a program, $d$ a call depth, $\sigma$ an initial
state, $p$ a procedure index with body $body$, $args$ a list of argument expressions, and
$x$ the local variable that receives the result. Suppose each argument is safe in $\sigma$,
$d<\mathrm{lim.depth}$, and the body executes at depth $d+1$ from the state whose local
frame is initialized with the evaluated arguments and whose memory is $\sigma.mem$, within
time $T\prime$, ending in a state $\sigma\prime$ such that updating local $x$ of $\sigma$
to $\sigma\prime.loc(0)$ and using memory $\sigma\prime.mem$ satisfies $Q$. If the sum
of argument evaluation costs plus $2+T\prime$ is at most $T$, then calling $p$ with
$args$ from $\sigma$ ends within time $T$ in a state satisfying $Q$.
\end{theorem}

\begin{definition}[Variable expression]
\lean{Light.v}
\label{Light.v}
\leanok
\uses{Light.Expr}
For each natural-number variable index $x$, the corresponding variable expression is the
expression representing that variable.
\end{definition}

\begin{definition}[Natural-number constant expression]
\lean{Light.k}
\label{Light.k}
\leanok
\uses{Light.Expr}
For each natural number $n$, the corresponding constant expression is the expression
representing the constant $n$.
\end{definition}

\begin{definition}[Memory-load expression]
\lean{Light.M}
\label{Light.M}
\leanok
\uses{Light.Expr}
For an expression $a$, the memory-load expression associated with $a$ is the expression
that loads from $a$.
\end{definition}

\begin{definition}[Non-strict comparison of expressions]
\lean{Light.Cond.le}
\label{Light.Cond.le}
\leanok
\uses{Light.Cond, Light.Expr, Light.k}
For expressions $a$ and $b$, the condition $a \le b$ means that $a$ is strictly less
than the expression obtained by adding the constant $1$ to $b$.
\end{definition}

\begin{definition}[Conditional on expression inequality]
\lean{Light.Stmt.iteNe}
\label{Light.Stmt.iteNe}
\leanok
\uses{Light.Expr, Light.Stmt}
Given expressions $a,b$ and statements $s_1,s_2$, define the conditional statement that
executes $s_1$ when $a=b$ and $s_2$ otherwise.
\end{definition}

\begin{theorem}[Meaning of the less-than-or-equal condition]
\lean{Light.Cond.holds_le}
\label{Light.Cond.holds_le}
\leanok
\uses{Light.Cond.Holds, Light.Expr, Light.Expr.val, Light.State}
\difficulty{1}
For a state $\sigma$ and expressions $a,b$, the condition written $a\le' b$ holds in
$\sigma$ if and only if the value of $a$ in $\sigma$ is at most the value of $b$ in
$\sigma$.
\end{theorem}

\begin{theorem}[Natural cast addition under truncation]
\lean{Light.toNat_natCast_add_natCast}
\label{Light.toNat_natCast_add_natCast}
\leanok
\difficulty{3}
For natural numbers $a$ and $b$, the natural number obtained by converting their integer
sum to a natural number is $a+b$.
\end{theorem}

\begin{definition}[Standard resource limits]
\lean{Light.Std}
\label{Light.Std}
\leanok
\uses{Light.Limits}
A limit record with word bound $w$, space bound $s$, and depth bound $d$ is standard if
$s\le w$ and $100\le w$.
\end{definition}

\begin{theorem}[Absolute bound for an address]
\lean{Light.Limits.Addr.abs_le}
\label{Light.Limits.Addr.abs_le}
\leanok
\uses{Light.Limits.Addr}
\difficulty{3}
Let a limit specification have space bound $S$ and word bound $W$, and let $a$ be an
integer address satisfying $0\le a<S$. If $S\le W$, with the natural number $S$ viewed
as an integer, then $|a|\le W$.
\end{theorem}

\begin{theorem}[Address bound from a word limit]
\lean{Light.Limits.addr_of_lt}
\label{Light.Limits.addr_of_lt}
\leanok
\uses{Light.Limits.Addr, Light.Limits.Addr.abs_le}
\difficulty{4}
Let a machine configuration have word limit $w\in\mathbb Z$ and space limit $s\in\mathbb
N$. If $s\le w$ and $x<s$, then the integer address $x$ is valid, meaning $0\le x<s$,
and its absolute value satisfies $|x|\le w$.
\end{theorem}

\begin{theorem}[Absolute bound below the space limit]
\lean{Light.Limits.abs_le_of_le}
\label{Light.Limits.abs_le_of_le}
\leanok
\difficulty{3}
Let a limit specification have natural-number space bound $S$ and integer word bound
$W$. If $S\le W$, with $S$ viewed as an integer, and a natural number $x$ satisfies
$x\le S$, then $|(x)|\le W$, where $x$ is viewed as an integer inside the absolute
value.
\end{theorem}
